%% 几何 Satake 对应
\section[几何 Satake 对应]{几何 Satake 对应（5 月 4 日）}

\srcpage{74}

几何 Satake 对应是 Fourier 级数理论的一个非阿贝尔对应物。嗯，做这件事有不同方式，正如我们已经看到的。它是「$\Z$-分次向量空间等价于 $\G_m$-表示」这一事实的推广。上次我们看到，$\Z$-分次向量空间等价于 $\mathrm{GL}_1$ 的仿射 Grassmann 簇上的层。另一边，$\G_m$-表示可以解释为 $\mathrm{Loc}_{\G_m}(\P^1)$ 上的拟凝聚层。

我们将推广这一点，问对非阿贝尔的 $G$，$\Gr_G$ 上的层是什么。在非阿贝尔情形，双陪集更棘手：一般说来它们与单陪集不同。我们是靠考虑「意大利饺」上的主 $G$-丛到达仿射 Grassmann 簇的，想的是圆盘上带原点之外等同的丛。具体地，$LG_+\backslash\Gr_G=LG_+\backslash LG/LG_+$。与此相关，
\[
  LG_-\backslash LG/LG_+ \cong \Bun_G(\P^1).
\]\yizhu{这里 $LG=G(K_x)=G(\C((t)))$，$LG_+=G(\cO_x)=G(\C[[t]])$，而 $LG_-=G(\C[t^{-1}])$；下标中的正负号指 Laurent 级数的非负部分与非正部分。}{I（术语背景）}

$\Gr_G$ 上的 $LG_+$-轨道在「包含于某轨道的闭包中」这一序下构成一个偏序集。这些轨道都是有限维的，且每个轨道闭包只含有限多个轨道。

\begin{example}
设 $G=\mathrm{GL}_2$。下面是 $\P^1$ 上的一些秩二向量丛，它们的主 $\mathrm{GL}_2$-丛即其标架丛：$\cO\oplus\cO(1)$、$\cO\oplus\cO$、$\cO(2)\oplus\cO(2)$，等等；而 Grothendieck–Birkhoff 定理说，所有主 $\mathrm{GL}_2$-丛在同构意义下都这样出现。

丛的次数（同样地，$\int_{\P^1}c_1(P)$）定义了一个到 $\Z$ 的映射；$0$ 的原像恰好是主 $\mathrm{SL}_2$-丛；这是一个连通的紧空间。它包含形如 $\cO(n)\oplus\cO(-n)$ 的丛，而每个 $n$ 自成一个轨道。$\cO(n)\oplus\cO(-n)$ 的轨道闭包包含 $\cO(n+1)\oplus\cO(n+1)$ 的轨道：你可以把粘合函数在族中变动，它可能在某处出现零点而退化。所以偏序集是 $\mathbb{N}$，且所有轨道都是有限维的。
\end{example}

一般地，$LG_-=G(\C[t^{-1}])$。$\Gr_G$ 包含一个位于各胞腔中心的格 $\Gr_T$（TODO：这很怪，我可能听错了什么），它们正是 $\G_m$-作用的不动点。而 $LG_+\backslash\Gr_G$ 与（TODO：这对我也不太说得通）$\Lambda/W$ 等同。这就是余权格。抽象地说，这是一个权格；由半单群的分类，我们知道这在告诉我们某个别的约化群 $G^\vee$ 的表示论。

我们想把这一切漂亮地打包：找一个群 $G^\vee$，使它有极大环面 $T^\vee$ 与相同的 Weyl 群，并满足 $\Lambda=\Hom(T^\vee,\G_m)$。也就是说：轨道的集合看起来很像某个约化群的不可约表示，那就让我们去把这个群找出来。事实上，我们将从球面 Hecke 范畴 $\mathrm{Shv}_{LG_+}(\Gr_G)$ 构造 $G^\vee$。这里的「层」可以指 D-模，也可以指一种叫做反常层（perverse sheaf）的东西（某些在轨道上局部常值的层，即轨道上带平坦联络的向量丛；等价地，轨道上按轨道维数移位的局部系统）。\yizhu{反常层（perverse sheaf）并非「反常」，而是满足一套支撑与对偶条件的复形；它在几何表示论中充当「层论的相交上同调」。}{I（术语背景）}对这个层范畴，单对象与轨道上的不可约局部系统一一对应；而 $LG_+$ 轨道都是单连通的（它们都是 $G$ 上的仿射丛或部分旗簇），所以我们理解这些局部系统。一般地，层或 D-模的范畴很复杂，但这一个将是半单的！（若我们想让它成为 $G^\vee$ 的表示范畴，这本是意料之中——但它仍然是个惊喜。）

这是 Grassmann 簇几何特有的东西，与维数的奇偶性有关。结果是：在 $\Gr_G$ 的每个连通分支上，轨道的复维数模 2 是常值。这限制了轨道上的系统与其闭包之间可能的相互作用——多出来的东西是余维 2，这迫使某些 $\Ext^1$ 项消失。（完整证明稍复杂些，但基本想法就是这样。）

于是我们得到：球面 Hecke 范畴是半单阿贝尔范畴，但它还有更多：它是对称幺半的。物理告诉我们，这是因为它是 4 维量子场论中线算子的范畴；但还有进一步的理由。某种意义上这里有两个结构在起作用：因为这是双陪集空间 $LG_+\backslash LG/LG_+$ 上的层，总有一个幺半结构，虽然它不自动交换。但还有一个编织乘法，来自一个完全不同的故事：$\Gr_G$ 是二重环路空间，所以它上面的层获得 $E_2$-代数结构。

Beilinson–Drinfeld 找到了 $E_2$-结构这一概念的一个纯代数几何类比，即一个 adelic 版本，并证明 $\Gr_G$ 具有该结构。它同样关乎曲线上点的碰撞，但他们用所谓 BD Grassmann 簇来形式化，它并不真有乘积，却仍编码了乘法。我们固定一条曲线 $C$，比如 $\A^1$，并考虑一个纤维丛 $\Gr_{G,C^n}\to C^n$。

\srcpage{75}

对 $n=1$，$\Gr_{G,C}$ 在 $x$ 处的纤维是 $C$ 上带 $C\setminus x$ 上平凡化的主 $G$-丛空间。要精确地定义它，应当在族中做，或描述其点函子。

对 $n=2$，我们在曲线的两份拷贝上有族，所以 $(x,y)\in C^2$ 处的纤维是 $C$ 上带 $C\setminus\{x,y\}$ 上平凡化的主 $G$-丛空间。它有两个分层：$x\ne y$ 与 $x=y$。对 $x\ne y$，$\Gr_{G,C^2}|_{(x,y)}\cong \Gr_{G,C}|_x\times\Gr_{G,C}|_y$：在 $x,y$ 之外平凡化的丛，等价于「在 $x$ 处平凡化的丛」加「在 $y$ 处平凡化的丛」的数据。所以对角线之外我们就是一个积。而当 $x=y$ 时，$\Gr_{G,C^2}|_{(x,x)}\cong\Gr_{G,C}|_x$：没有额外数据。

这给了我们某种有趣的东西：我们把某个东西的两份拷贝粘成了它的一份拷贝。例如对 $G=\mathrm{GL}_1$，在 $C^2$ 中对角线之外这是 $\Z\times\Z$，而对角线上只是 $\Z$。这很怪：一般地，代数几何中维数在特殊化时可能变大；这里在对角线之外反而大了一倍！解决办法是：你需要有限维性来证明那个事实，而这里一切都是无穷维的。

一般地，$\Gr_{G,C^n}$ 有两个关键结构。第一，它在 $(x_1,\dots,x_n)$ 处的纤维只依赖于子集 $\{x_1,\dots,x_n\}$ 本身，而不依赖于次序。这意味着当你限制到对角子空间时，得到预期的 $\Gr_{G,C^k}$（$k<n$）。于是我们可以（TODO：我想）把这些粘起来，取余极限，并把它定义在 Ran 空间 $\mathrm{Ran}(C)$ 上，即 $C$ 的有限子集构成的空间。你也可以把它定义为 $\mathrm{colim}_{I\to J}C^I$。作为一个空间这有点怪，但常常方便地把它想成一个滤过系统：Ran 空间上的层，就是对每个 $C^I$ 给一个层，并与 $I\to J$ 的拉回相容。

第二个重要事实是因子化：这个构造对不交并是乘性的，即在 $\{x_i\}\amalg\{y_i\}$ 处的纤维，是在 $\{x_i\}$ 与 $\{y_i\}$ 处纤维的积。这是结构而非数据，是 $E_2$-代数结构的一个版本，但更代数几何化。它并不给出像小球盘空间那样的真正乘法映射；但因为有一些办法把 Grassmann 簇的两份拷贝粘成一份，我们就有某种像是族中的乘法。这是因子化空间的一个例子，而 $\Gr_G$ 有这个结构；它是编码那个有点不寻常的代数结构的一种方式，该结构在同伦范畴中就是 $E_2$-结构。在代数几何中我们没有这个 $E_2$-结构，但我们有时确实得到乘法：它定义了 Grassmann 簇上那些可以取极限的对象上的乘法。

例如，设 $\cF$ 与 $\cG$ 是 $\Gr_G$ 上的层。我们将用算子乘积展开把它们乘起来：把 $\cF$ 想在点 $x$、$\cG$ 想在点 $y$。于是我们在 $\{x,y\}$ 上有 $\cF\boxtimes\cG$，并且可以在变动 $x$、让它靠近 $y$ 时这样做。我们想把它延拓到对角线上，而这里取决于我们有什么样的层：对凝聚层有问题，但对更拓扑的东西，有办法做到，那就是用邻近轮（nearby cycles），它是极限的拓扑版本。于是我们可以用它把 $\cF*\cG$ 定义为 $\cF\boxtimes\cG$ 在对角线上的邻近轮。共同的主题是：接近拓扑的东西总有某种极限／延拓。

这并不总是正确的做法——在这个乘法下，$\Gr_G$ 上的层看起来像编织幺半结构，但它不结合！问题在于我们需要对几何施加更多约束：如果你限制到 $LG_+$-等变层，那就行了，给出 $\mathrm{Sph}$ 上结合的编织幺半结构。（缺少的成分是一个紧性假设——某种意义上 Grassmann 簇有点像 $E_2$-空间，我们把它上面的层想成有 $E_2$-结构，但这实际上只对足够小的层子范畴成立，例如 $LG_+$-等变的那些。）

你可以用两种方式把 $\Gr_G$ 上的两个层叠起来：可以竖直地叠，这就是 $LG_+\backslash LG/LG_+$ 的 Hecke 运算，它是结合的；还有因子化运算，让点碰撞，它是编织的。这两个乘法都结合且彼此交换，因此 Eckmann–Hilton 告诉我们得到一个对称幺半结构，并以代数几何方式编码。我们把三维的对称性拆成了 $2+1$：平面内两个方向（曲线），以及离开它的一个方向（叠／Hecke 运算）。

因子化这一概念——把类似 $E_2$-结构的东西代数几何化——是几何 Langlands 对应产生的最重要的东西之一。它出自物理，那是特征零；但现在我们可以把同样的因子化概念搬到有限域上的曲线，或（如 Scholze 所做的）搬到 $p$ 进域上工作。

我们将通过 Tannakian 重构得到 $G^\vee$。我们已经造出一个对称幺半范畴 $(\cC,\otimes)$，于是得到一个叠 $\Spec(\cC,\otimes)$；对足够好的 $\cC$，我们可以把 $\cC$ 还原为 $\QCoh(\Spec(\cC,\otimes))$。

\srcpage{76}

作为一种特殊情形——更熟悉的 Tannakian 重构概念——我们可以指定一个纤维函子，即到 $(\Vect_k,\otimes)$ 的忠实正合对称幺半函子。这个额外数据告诉我们 $\Spec(\cC,\otimes)$ 有一个 $k$-点，即一个从 $\Spec(\Vect_k,\otimes)=\Spec k$ 出发的映射，且这个映射是平坦覆盖。因此 $\Spec(\cC,\otimes)\cong \mathrm{pt}/BG^\vee$，对某个群 $G^\vee$，称为 $(\cC,\otimes)$ 的 Galois 群或 Tannakian 群；而 $G^\vee$ 是纤维函子的自同构群。

在我们的情形，$\cC$ 是 $\Gr_G$ 上层的范畴，我们有好用的对称幺半结构。接下来我们需要一个纤维函子，即到 $\Vect_k$ 的遗忘函子。换句话说，一个层的「底层向量空间」是什么？有什么好办法来度量它，既忠实又对称幺半？有各种想法；最短的说法是：取上同调。但那可能不是最有洞见的说法；我们将换一种方式想纤维函子，造一个映射
\[
  \mathrm{Shv}_{LG_+}(\Gr_G)\longrightarrow \mathrm{Shv}(\Gr_T).
\]
由于 $\Gr_T$ 是一个格，我们得到一个到 $\Rep(T^\vee)$ 的函子，其中 $T^\vee:=\Hom(\Lambda,\G_m)$。现在 $\Lambda\subset\Gr_G$，所以我们可以把层限制到一个离散集上。限制之后如何度量是有选择的，而你随手写下的朴素函子确实是错的；但这样想仍然更容易、也更清楚。

下次我们会更详细地讨论答案（抛物限制），但数据的快速总结是：在 $G$ 与 $T$ 之间有一个 Borel 子群 $B$，例如 $G=\mathrm{GL}_n$ 时是上三角矩阵。这给我们一个对应，因为 $T$ 不只是 $G$ 的子群，还是 $B$ 的商：
\[
  \begin{tikzcd}
    & \Gr_B \arrow[dl] \arrow[dr] & \\
    \Gr_G & & \Gr_T\cong\Lambda.
  \end{tikzcd}
\]
那么我们的纤维函子将是（还能是什么？）层上的推拉算子，给出 $\Lambda$ 上的一个层，等价地，$T^\vee$ 的一个表示。我们的纤维函子遗忘这个作用，但记住它给了我们额外数据：球面 Hecke 范畴有一个从 $\mathrm{pt}/G^\vee$ 出发的映射。
